\begin{abstract}
Singapore requires low-carbon technologies that can deliver more than 0.25 MtCO$_2$e/y of abatement at less than S\$100/tCO$_2$e. This study evaluates a literature-anchored nuclear-assisted steam-methane-reforming (SMR) plant with carbon capture and storage (CCS). The final process basis is INL's NGNP high-temperature case: an 871$^\circ$C reformer at steam/carbon 3.0, supplied with 900$^\circ$C nuclear process heat from a 925$^\circ$C HTGR outlet, with 78.1\% methane conversion, 88\% PSA recovery and carbon capture~\cite{inltev953,inltev961}. JAEA GTHTR300C evidence supplies the 600 MWth module, helium/IHX architecture and one-module cogeneration economics~\cite{jaeri2004gthtr300c,nishihara2007potential}. At 85\% availability the 130 MMSCFD source service corresponds to approximately 97,946 tH$_2$/y. Extending the source balance with upstream-gas, nuclear-lifecycle, CCS-transport and Singapore economic accounting gives approximately 0.917 MtCO$_2$e/y lifecycle abatement. The conservative zero-value-cogeneration case gives approximately S\$30.2/tCO$_2$e, satisfying the CN4252 cost threshold without electricity revenue; the mature S\$150/MWh screening case gives approximately $-$S\$194.7/tCO$_2$e because modelled annual candidate cost falls below the baseline under that electricity-value assumption. The result is a \textbf{conditional model result}, not observed Singapore commercial feasibility: it depends on mature HTGR design-study economics, an enlarged high-temperature IHX/secondary-loop treatment, cross-border CO$_2$ storage and lifecycle screening proxies.
\end{abstract}
